\documentclass[11pt]{article}
\usepackage[margin=0.85in]{geometry}
\usepackage{amsmath,amssymb,amsthm}
\usepackage{mathtools}
\usepackage{booktabs}
\usepackage[dvipsnames]{xcolor}
\usepackage{hyperref}
\hypersetup{colorlinks=true,linkcolor=blue!50!black,urlcolor=blue!50!black}

\newtheorem{lemma}{Lemma}
\newtheorem{theorem}{Theorem}

\newcommand{\tphi}{\mathrm{tr}(\phi^3)}
\newcommand{\Zr}{\mathcal{Z}}
\newcommand{\Fr}{F}

\title{Locality and unitarity emerge from hidden zeros in $\tphi$:\\
a by-hand proof at four, five, and six points}
\author{Jony}
\date{April 25, 2026}

\begin{document}
\maketitle

\begin{abstract}
We prove that the cyclic hidden one-zeros of color-ordered
$\tphi$ tree amplitudes force every non-local coefficient of the
general rational ansatz to vanish, and force every triangulation
coefficient to take a single common value, at $n=4,5,6$ external
legs.  The proof uses one move applied at each cyclic zone:
\emph{send a chosen subset of free chords to infinity and read off
coefficients of the surviving rational function in the remaining
variables}.  No Laurent expansions are required.  Locality (the
survivors are exactly triangulations) and unitarity (they share a
common scalar) emerge together.
\end{abstract}

% ==============================================================
\section{Setup}
% ==============================================================

Label the legs of an $n$-point amplitude cyclically by $1,\dots,n$
and let $X_{ij}=(p_i+\cdots+p_{j-1})^2$ for $1\le i<j\le n$, $j-i\ge 2$,
$(i,j)\ne(1,n)$.  There are $d=n(n-3)/2$ such planar Mandelstam
chords.  A \emph{triangulation} is a maximal non-crossing set of
$n-3$ chords, and the tree amplitude is
$A_n = \sum_{T}\prod_{c\in T}1/X_c$ summed over triangulations.
We consider the most general rational ansatz with mass dimension
$-2(n-3)$ and only planar simple poles,
\begin{equation}\label{eq:ansatz}
B \;=\; \sum_{w}\frac{a_w}{\prod_{c\in w}X_c},
\end{equation}
where $w$ runs over size-$(n-3)$ multisets of chord indices.  A
multiset is \emph{local} if its chords form a triangulation.

Non-planar invariants $c_{a,b}=2p_a\cdot p_b$ satisfy the mesh identity
$c_{a,b}=X_{a,b+1}+X_{a+1,b}-X_{a,b}-X_{a+1,b+1}$ (with $X_{a,a+1}=0$).
The hidden one-zero $\Zr_r$ is the linear subspace
$c_{r,r+2}=\cdots=c_{r,r+n-2}=0$, $n-3$ constraints.  Solving these
via the mesh identity gives the master substitution at $\Zr_r$:
every chord $X_{r+1,j}$ becomes a $\mathbb{Q}$-linear function of
the remaining (\emph{free}) chords, with the special chord
$X_*:=X_{r,r+2}$ appearing on every right-hand side.  We write
$\Fr_r$ for the free set excluding $X_*$.  Rodina has shown
$A_n|_{\Zr_r}\equiv0$; our goal is the converse: assuming
$B|_{\Zr_r}=0$ for all $r$, show $B=a\cdot A_n$ for some scalar $a$.

\begin{lemma}[The kill move]\label{lem:move}
Suppose $B|_{\Zr_r}=0$.  Let $S\subseteq \Fr_r\cup\{X_*\}$ be any
subset of free chords held arbitrary, and let every chord outside
$S$ go to $\infty$.  Whenever a non-bare master substitution
$X_s = X_f - X_*$ has its companion $X_f$ outside $S$, also let
$X_f \to \infty$ (otherwise $1/X_s = 1/(X_f-X_*)$ would survive as
a non-monomial rational factor, blocking coefficient extraction).
Then the surviving terms of $B|_{\Zr_r}$ form a Laurent polynomial
in $S$, and each coefficient on a distinct monomial vanishes.
\end{lemma}
\begin{proof}
After the master substitution, $B|_{\Zr_r}$ is a rational function
in the $|\Fr_r|+1$ free chords.  Sending a chord to infinity drops
every monomial of $B$ whose denominator contains it.  The
companion condition ensures all surviving monomials are Laurent
monomials in $S$; distinct multisets give distinct monomials, so
linear independence forces individual vanishing.
\end{proof}

\noindent We use Lemma~\ref{lem:move} in two flavors:
\begin{itemize}\setlength\itemsep{0pt}
  \item \textbf{Step 1} drops $X_*$ from $S$.  Surviving multisets
  are subsets of $\Fr_r$, and coefficients are forced to zero
  individually.
  \item \textbf{Step 2} keeps $X_*\in S$ and exploits
  $X_{\text{bare}}=-X_*$.  All non-bare substituted chords get
  dropped (by sending their companions outside $S$ to infinity).
  The surviving polynomial in $S$ pairs multisets glued by the
  sign of $-X_*$, producing equalities of original coefficients.
\end{itemize}

% ==============================================================
\section{Four points}
% ==============================================================

The square has only $X_{13}, X_{24}$ and the ansatz
$B = a_{X_{13}}/X_{13}+a_{X_{24}}/X_{24}$ has \emph{no} non-local
multisets: with $n-3=1$, every single chord is itself a
triangulation, and there is no room for a crossing.  Locality is
automatic.  The single zero $\Zr_1$ gives $X_{24}=-X_{13}$, so
$B|_{\Zr_1}=(a_{X_{13}}-a_{X_{24}})/X_{13}=0$, hence
$a_{X_{13}}=a_{X_{24}}$.  This is unitarity.  $B = a\cdot A_4$.

% ==============================================================
\section{Five points}
% ==============================================================

\subsection{Setup}

The pentagon has $5$ chords, all class $1$ (short).  The ansatz
has $\binom{5+1}{2}=15$ coefficients: $5$ tree triples (the fan
triangulations), $5$ double poles, $5$ crossing pairs.  We must
kill the latter $10$ and equate the $5$ tree triples.  The five
zones are
\[
\begin{array}{r|ll|c|c}
\text{Zone} & \multicolumn{2}{c|}{\text{Master subs}} & X_* & \Fr_r\\\hline
\Zr_1 & X_{24}=X_{14}-X_{13} & X_{25}=-X_{13} & X_{13} & \{X_{14},X_{35}\}\\
\Zr_2 & X_{35}=X_{25}-X_{24} & X_{13}=-X_{24} & X_{24} & \{X_{25},X_{14}\}\\
\Zr_3 & X_{14}=X_{13}-X_{35} & X_{24}=-X_{35} & X_{35} & \{X_{13},X_{25}\}\\
\Zr_4 & X_{25}=X_{24}-X_{14} & X_{35}=-X_{14} & X_{14} & \{X_{24},X_{13}\}\\
\Zr_5 & X_{13}=X_{35}-X_{25} & X_{14}=-X_{25} & X_{25} & \{X_{35},X_{24}\}\\
\end{array}
\]
Each zone has one special $X_*$, one non-bare substitution $X_s = X_f - X_*$
with companion $X_f\in\Fr_r$, one bare $X_{\text{bare}}=-X_*$, and two free chords.

\subsection{Step 1 at $\Zr_1$ (three kills)}

Send $X_{13}\to\infty$ holding $X_{14},X_{35}$ fixed: then
$X_{24},X_{25}\to-\infty$ and every $B$-monomial involving any of
$X_{13},X_{24},X_{25}$ drops out.  Surviving (with $a_{ij,kl} := a_{X_{ij}X_{kl}}$):
\[
\sum_{w\subseteq\{X_{14},X_{35}\}} \frac{a_w}{\prod X_w}
= \frac{a_{14,14}}{X_{14}^2} + \frac{a_{14,35}}{X_{14}X_{35}} + \frac{a_{35,35}}{X_{35}^2}=0,
\]
so by independence $a_{14,14}=a_{14,35}=a_{35,35}=0$.

\subsection{Step 2 at $\Zr_1$ (pair equalities)}

To use the bare relation $X_{25}=-X_{13}$ cleanly we must drop the
non-bare substituted chord $X_{24}=X_{14}-X_{13}$.  Sending
$X_{14}\to\infty$ does this: then $X_{24}\to\infty$ as well, and
\emph{both} $X_{14}$- and $X_{24}$-involving monomials drop from
the surviving sum.  This is essential --- if we kept $X_{14}$
finite while letting $X_{24}$ contribute, every $X_{24}$-monomial
would carry a factor $1/X_{24} = 1/(X_{14}-X_{13})$, which is not
a Laurent monomial in $\{X_{13}, X_{14}\}$ and would block clean
coefficient extraction.

Holding $X_{13},X_{35}$ fixed and sending $X_{14}\to\infty$, the
survivors are multisets in $\{X_{13},X_{25},X_{35}\}$, glued by
$X_{25}=-X_{13}$.  Below we abbreviate $a_{X_{ij}X_{kl}}$ as $a_{ij,kl}$:
\[
\frac{a_{13,13}}{X_{13}^2} - \frac{a_{13,25}}{X_{13}^2}
+\frac{a_{25,25}}{X_{13}^2}
+ \frac{a_{13,35}}{X_{13}X_{35}} - \frac{a_{25,35}}{X_{13}X_{35}}
+ \frac{a_{35,35}}{X_{35}^2} = 0.
\]
Collecting on $1/X_{13}^2,\ 1/(X_{13}X_{35}),\ 1/X_{35}^2$:
\[
\boxed{a_{13,13}-a_{13,25}+a_{25,25}=0,\qquad a_{13,35}=a_{25,35},\qquad a_{35,35}=0\ \text{(already from Step 1).}}
\]
The first identity, combined with Step 1 kills of $a_{13,13}$ (via
$\Zr_3,\Zr_4$) and $a_{25,25}$ (via $\Zr_2,\Zr_3$), forces the
crossing pair $a_{13,25}=0$ --- a second route to the same kill.
The middle identity equates two fan triangulations sharing $X_{35}$.

\subsection{Cycling and the union}

Cycling $1\to 2\to\cdots\to 5$ rotates Step 1 and Step 2 across all five zones.  Below we abbreviate $a_{X_{ij}X_{kl}}$ as $a_{ij,kl}$:
\begin{center}\small
\begin{tabular}{c|c|c}
Zone & Step 1 kills & Step 2 unitarity equality\\\hline
$\Zr_1$ & $a_{14,14},\ a_{35,35},\ a_{14,35}$ & $a_{13,35}=a_{25,35}$\\
$\Zr_2$ & $a_{25,25},\ a_{14,14},\ a_{14,25}$ & $a_{14,24}=a_{13,14}$\\
$\Zr_3$ & $a_{13,13},\ a_{25,25},\ a_{13,25}$ & $a_{25,35}=a_{24,25}$\\
$\Zr_4$ & $a_{24,24},\ a_{13,13},\ a_{13,24}$ & $a_{13,14}=a_{13,35}$\\
$\Zr_5$ & $a_{35,35},\ a_{24,24},\ a_{24,35}$ & $a_{24,25}=a_{14,24}$\\
\end{tabular}
\end{center}

\paragraph{Locality.}
The Step 1 union covers the 5 double poles ($a_{ij,ij}$, each in
two zones) and the 5 crossing pairs ($a_{14,35}$ at $\Zr_1$,
$a_{14,25}$ at $\Zr_2$, $a_{13,25}$ at $\Zr_3$, $a_{13,24}$ at
$\Zr_4$, $a_{24,35}$ at $\Zr_5$).  All ten non-locals are killed.

\paragraph{Unitarity.}
The five Step 2 equalities form a $5$-cycle on the five fan triangulations:
\(
a_{13,35} \overset{\Zr_1}{=} a_{25,35}
          \overset{\Zr_3}{=} a_{24,25}
          \overset{\Zr_5}{=} a_{14,24}
          \overset{\Zr_2}{=} a_{13,14}
          \overset{\Zr_4}{=} a_{13,35}.
\)
All five fan coefficients are equal to a common $a$, hence
$B = a\cdot A_5$.

\begin{theorem}[$n=5$]\label{thm:n5}
The five hidden one-zeros of the pentagon force $B = a\cdot A_5$
for a unique scalar $a$.  Locality and unitarity emerge
simultaneously from the kill-and-equate enumeration.
\end{theorem}

% ==============================================================
\section{Six points}
% ==============================================================

\subsection{Setup}

The hexagon has $9$ chords: six short ($X_{13},X_{24},X_{35},X_{46},X_{15},X_{26}$,
each enclosing one vertex) and three long ($X_{14},X_{25},X_{36}$,
each enclosing two).  The ansatz has $\binom{9+2}{3}=165$
coefficients, of which $14$ are tree triples ($C_4=14$
triangulations) and $151$ are non-local.  Each $\Zr_r$ has three
master substitutions, two non-bare and one bare.  For $\Zr_1$:
\[
X_{24}=X_{14}-X_{13},\qquad X_{25}=X_{15}-X_{13},\qquad X_{26}=-X_{13},
\qquad X_*=X_{13}.
\]
The free set is
$\Fr_1 = \{X_{35},X_{46},X_{15},X_{14},X_{36}\}$ (3 short, 2 long).
The other zones are cyclic rotations.

\subsection{Step 1 at $\Zr_1$ (35 kills)}

Send $X_{13}\to\infty$, $\Fr_1$ fixed.  Then $X_{24},X_{25},X_{26}\to\pm\infty$,
and the surviving terms are exactly the multisets in $\Fr_1$.
Writing $a_{abc} := a_{X_aX_bX_c}$:
\[
\sum_{(a,b,c)\subseteq\Fr_1}\frac{a_{abc}}{X_aX_bX_c}=0
\quad\Longrightarrow\quad
\boxed{a_{abc}=0\ \text{for every multiset of size 3 from }\Fr_1.}
\]
There are $\binom{5+2}{3}=35$ such multisets.

\subsection{Step 2 at $\Zr_1$ (equalities and 3-term identities)}

To use the bare relation $X_{26}=-X_{13}$ cleanly we must drop
\emph{both} non-bare substitutions $X_{24}=X_{14}-X_{13}$ and
$X_{25}=X_{15}-X_{13}$.  If either survived in the sum, its
$1/X_{24}=1/(X_{14}-X_{13})$ or $1/X_{25}=1/(X_{15}-X_{13})$ factor
would not be a Laurent monomial in our chosen variables, blocking
clean coefficient extraction.  Sending both companions
$X_{14},X_{15}\to\infty$ achieves this: $X_{24},X_{25}\to\infty$
also, and every $B$-monomial involving any of $X_{14},X_{24},X_{15},X_{25}$ drops.

Keep $S=\{X_{13},X_{35},X_{46},X_{36}\}$ arbitrary.  Surviving
multisets draw from $\{X_{13},X_{26},X_{35},X_{46},X_{36}\}$, glued
by $X_{26}=-X_{13}$.  Each surviving $B$-monomial becomes a
Laurent monomial in $S$, with sign $(-1)^k$ where $k$ counts
factors of $X_{26}$.  Writing $a_{ab,c}:=a_{X_aX_bX_c}$:
\[
\boxed{a_{13^3}-a_{13^2,26}+a_{13,26^2}-a_{26^3}=0,\quad
a_{13^2,b}-a_{13,26,b}+a_{26^2,b}=0,\quad
a_{13,a,b}=a_{26,a,b}}
\]
for $X_a, X_b\in\{X_{35},X_{46},X_{36}\}$.  The third identity is
a direct equality of two coefficients; whether each is between two
triangulations or two non-locals depends on $\{X_a, X_b\}$ and is
sorted out by the orbit accounting below.  Cyclic rotations to
$\Zr_2,\dots,\Zr_6$ give the full Step 2 family.

\subsection{Cycling: locality and unitarity}

\paragraph{Locality at $n=6$.}  Step 1 across all six zones kills
every multiset whose support sits inside some $\Fr_r$ ($129$ of
the $151$ non-local multisets, after accounting for overlap).  The
remaining $22$ non-local multisets touch every vertex; they fall
into five cyclic-symmetry orbits:
\[
\text{(a) } X_{r,r+2}^2 X_{r+3,r+5},\;
\text{(b) } \{X_{r,r+2}, X_{r+1,r+3}, X_{r+3,r+5}\},\;
\text{(c) } \{X_{r,r+2}, X_{r+1,r+3}, X_{r+4,r+6}\},\;
\]
\[
\text{(d) } \{X_{r,r+2}, X_{r+3,r+5}, X_{r,r+3}\},\;
\text{(e) } \{X_{14}, X_{25}, X_{36}\}.
\]
For each orbit, two cyclically-shifted Step 2 identities (case A
at $\Zr_r$, case B at $\Zr_{r-1}$) give both an equality
$a_M = a_{M'}$ \emph{and} a difference $a_M - a_{M'} = 0$, forcing
$a_M = 0$.  The detailed orbit verification is in the companion
note (\#17 page 4 et seq.); each verification is one application
of Step 2, so all $22$ remaining non-locals die, completing
locality.

\paragraph{Unitarity at $n=6$.}  The Step 2 identities
$a_{13,a,b}=a_{26,a,b}$ equate adjacent triangulation coefficients;
each pair $(X_{13},X_{26})$ corresponds to a flip in the hexagon's
associahedron $K_5$.  Cyclically, every flip in the
$14$-triangulation flip-graph is realized by some Step 2 equality,
the flip-graph is connected, so all $14$ tree-triple coefficients
equal a common $a$, and $B=a\cdot A_6$.

\begin{theorem}[$n=6$]\label{thm:n6}
The six hidden one-zeros of the hexagon force $B=a\cdot A_6$.  All
$151$ non-local coefficients vanish ($129$ by Step 1, $22$ by
cyclic-shifted Step 2 on five orbits) and the $14$ triangulation
coefficients collapse to a single common value via Step 2 flips.
\end{theorem}

\paragraph{General-$n$ outlook.}  The two-step move applies at every
$n$.  Open: (i) does cyclic-shifted Step 2 cover all multisets
outside any $\Fr_r$? (verified $n\le 6$, numerical $n=7$);
(ii) does Step 2 cover all flips of $K_{n-1}$? (verified $n\le 6$).
Closing both is the central goal of the ongoing program.

\end{document}
